\recipeTitle{Crescentine}

\recipeSubtitle{Le crescentine (che non si chiamano assolutamente "gnocco fritto") sono un piatto tipico dell'Emilia Romagna, in particolare della zona di Bologna. Si tratta di strisce di impasto lievitato fritte in strutto bollente, perfette per aperitivi o feste, da accompagnare con salumi, formaggi o marmellate.}

\recipeAuthor{Simone Balducci}

\recipeStory{I cräṡäntén i èn al magnèr dla fèsta a Bulåggna: ch’al sia un cumpléann, una laurea, un aperitévv, o una sèmpṡa znéṅna tra amîg, par stèr in cumpagné a n gh’è gnint ed mèi d’na cräṡänténa cun la mortadèla.}

% Put the image in recipes/nome-ricetta/images/ and update the path below.
% \recipePhoto{recipes/nome-ricetta/images/foto.jpg}

\recipeInfo
  {30 minuti}
  {20 minuti}
  {4 persone}
  {Facile}
  {Antipasto}

\begin{ingredients}
  \ingredient{Strutto di maiale}{20\,g}
  \ingredient{Farina 00}{500\,g}
  \ingredient{Latte intero}{250\,ml a temperatura ambiente}
  \ingredient{Lievito}{25\,g di lievito di birra fresco o 7\,g di lievito secco}
  \ingredient{Sale}{qb}
\end{ingredients}

\begin{tools}
  \tool{Ciotola}
  \tool{Tazza}
  \tool{Mattarello}
  \tool{Pentola per friggere}
\end{tools}

\begin{procedure}
  \step{Iniziare versando il latte a temperatura ambiente in una tazza e aggiungere il lievito. Utilizzando una forchetta, mescolare energeticamente fino a quando il lievito si sarà completamente sciolto.}
  \step{In una ciotola, riporre la farina, lo strutto e il sale. Scogliere lo strutto con le mani per integrarlo con la farina. A questo punto, versare il latte nella ciotola e iniziare a impastare.}
  \step{Una volta versato tutto il latte e ottenuto un impasto compatto, trasferirsi su una superficie piana leggermente infarinata e lavorare l'impasto fino ad ottenere un composto liscio ed elastico (approssimativamente 10/15 minuti).}
  \step{Lasciar lievitare l'impasto in una ciotola coperta con un canovaccio umido o pellicola trasparente, in un luogo caldo, fino a quando non raddoppia di volume (circa 1.5/2 ore).}
  \step{Stendere l'impasto con un mattarello fino a raggiungere uno spessore uniforme di circa 5\,mm.}
  \step{Tagliare l'impasto in strisce grandi a piacere.}
  \step{Friggere in abbondante strutto bollente fino a quando le crescentine non saranno gonfie e dorate, girandole a metà cottura per assicurare una doratura uniforme.}
  \stem{Servire con abbondanti salumi, formaggi o marmellate a piacere.}
\end{procedure}
